% =====================================================================
%  Lecture 08 -- Random Access Channel (RACH)             (English edition)
% =====================================================================
\documentclass[11pt,a4paper]{article}
\providecommand\ACNroot{../..}
\input{\ACNroot/templates/lecture-en}

\ACNset{number}{08}
\ACNset{title}{Random Access Channel (RACH)}
\ACNset{shorttitle}{Random Access (RACH)}
\ACNset{titleother}{کانال دسترسی تصادفی (RACH)}
\ACNset{part}{Part II -- 4G access and core (Lectures 5--13)}
\ACNset{time}{6--8 hours (3 h reading, 3--5 h exercises)}
\ACNset{level}{Intermediate}
\ACNset{prereq}{Lectures 6--7; slotted ALOHA; complex exponentials and correlation}
\ACNset{gen}{4}

\begin{document}
\ACNcover

% ---------------------------------------------------------------------
\section{Learning objectives}
After studying this lecture you should be able to:
\begin{objectives}
  \item \textbf{List} the situations in which an LTE UE must perform random access and
        \textbf{explain} why uplink synchronisation is the core problem.
  \item \textbf{Describe} the PRACH: preamble formats, cyclic prefix, guard time and the
        maximum cell range they imply.
  \item \textbf{Explain} Zadoff--Chu sequences, their correlation properties and how 64
        preambles are generated per cell with cyclic shifts.
  \item \textbf{Trace} contention-based (4-step) and contention-free random access, and the NR
        2-step variant.
  \item \textbf{Compute} timing advance values and the cell radius supported by a cyclic-shift
        configuration.
  \item \textbf{Model} preamble collisions and RACH overload, and \textbf{evaluate} countermeasures
        such as access class barring and backoff.
\end{objectives}

% ---------------------------------------------------------------------
\section{Prerequisite knowledge}
You need the LTE frame structure and uplink scheduling of Lecture~7 (in particular that PUSCH
transmissions of different UEs must arrive time-aligned at the eNB), the timing advance concept
of GSM (Lecture~1), the slotted ALOHA model of random access from introductory networking, and
the definition of the periodic cross-correlation of two complex sequences.

\paragraph{Key terms} The vocabulary of this lecture with its Persian equivalents:
\begin{keyterms}
Random access (RA) & \fa{دسترسی تصادفی} & Uncoordinated first uplink access of a UE. \\
PRACH & \fa{کانال فیزیکی دسترسی تصادفی} & Time--frequency resources reserved for preambles. \\
Preamble & \fa{پیش‌نوا} & Known sequence sent by the UE in Msg1. \\
Zadoff--Chu (ZC) sequence & \fa{دنبالهٔ زادوف--چو} & Constant-amplitude sequence with ideal cyclic autocorrelation. \\
Cyclic shift ($N_{\mathrm{CS}}$) & \fa{جابه‌جایی چرخشی} & Shift that creates orthogonal preambles from one root. \\
Timing advance (TA) & \fa{پیش‌فرستی زمانی} & Uplink advance compensating the round-trip delay. \\
RAR & \fa{پاسخ دسترسی تصادفی} & Msg2: TA, uplink grant and temporary identity. \\
Contention resolution & \fa{حل رقابت} & Msg4: tells which of the colliding UEs won. \\
Backoff & \fa{عقب‌نشینی} & Random wait before a new attempt. \\
Access class barring (ACB) & \fa{مسدودسازی ردهٔ دسترسی} & Probabilistic admission control before RA. \\
\end{keyterms}

% ---------------------------------------------------------------------
\section{Introduction and motivation}
Scheduled access needs two things the network does not have for a newly arriving UE: an
\emph{identity} to address it on the PDCCH and \emph{uplink time alignment}. A UE that has just woken
up knows the downlink timing (from the synchronisation signals) but not its distance to the eNB.
If it simply transmitted on PUSCH, its symbols would arrive late by the round-trip delay and destroy
the orthogonality of SC-FDMA with other users.

Random access is the procedure that solves both problems at once. The UE sends a short, robust
signal (a \emph{preamble}) in resources reserved for this purpose; the eNB measures its arrival time,
sends back a timing correction and a first grant, and the two sides then resolve the possibility that
several UEs chose the same preamble. The procedure is used surprisingly often -- at every call set-up
from idle, at every handover, after radio link failure and whenever an out-of-sync UE has data -- so its
capacity, latency and robustness matter. With massive IoT and 5G beam-based access, random access has
become one of the most actively engineered parts of the radio interface.

% ---------------------------------------------------------------------
\section{Fundamental concepts}

\subsection{When random access is triggered}
In LTE the MAC performs random access \cite{ts36321} for:
\begin{itemize}
  \item initial access from RRC\_IDLE (RRC connection establishment);
  \item RRC connection re-establishment after radio link failure;
  \item handover, to synchronise with the target cell;
  \item downlink data arrival for a connected UE that has lost uplink synchronisation;
  \item uplink data arrival when the UE is out of sync or has no scheduling-request resources;
  \item positioning, where the timing advance is needed.
\end{itemize}

\subsection{Contention-based versus contention-free}
In \textbf{contention-based} random access (CBRA) the UE picks a preamble at random from a common set;
two UEs may pick the same one, so a contention-resolution step is needed. In \textbf{contention-free}
random access (CFRA) the eNB assigns a dedicated preamble in advance (for handover or downlink data
arrival); no collision can occur and the procedure ends after the RAR.

\subsection{Zadoff--Chu sequences}
\begin{definitionbox}{Zadoff--Chu sequence}
For an odd prime length $N_{\mathrm{ZC}}$ and a root index $u\in\{1,\dots,N_{\mathrm{ZC}}-1\}$,
\begin{equation}\label{eq:zc}
  x_u(n) = \exp\!\left(-\jj\,\frac{\pi\,u\,n\,(n+1)}{N_{\mathrm{ZC}}}\right), \qquad n = 0,1,\dots,N_{\mathrm{ZC}}-1 .
\end{equation}
Its amplitude is constant, its periodic autocorrelation is zero for every non-zero shift, and the
cross-correlation of two different roots has constant magnitude $1/\sqrt{N_{\mathrm{ZC}}}$ (normalised).
\end{definitionbox}
These properties make ZC sequences ideal preambles: constant amplitude means a low peak-to-average
power ratio; zero autocorrelation means that \emph{cyclic shifts} of one root are orthogonal and the
delay of a received preamble shows up as a clean correlation peak. LTE uses $N_{\mathrm{ZC}} = 839$ for
preamble formats 0--3 \cite{ts36211}; the script \texttt{scripts/validate\_examples.py} verifies the
properties numerically for $u=25$ and $u=29$.

% ---------------------------------------------------------------------
\section{Detailed technical explanation}

\subsection{The PRACH}
The PRACH occupies six resource blocks (1.08\,MHz) in the uplink, in subframes given by the
\emph{PRACH configuration index} broadcast in SIB2. Inside this band the preamble uses a subcarrier
spacing of 1.25\,kHz, so its 839 subcarriers fit in about 1.05\,MHz. A preamble consists of a cyclic
prefix, the ZC sequence (800\,$\mu$s, or 1600\,$\mu$s repeated) and an implicit guard time:
\begin{table}[htbp]
\centering\small
\caption{LTE FDD preamble formats \cite{ts36211}. The guard time bounds the round-trip delay, hence the
cell range $d_{\max} = c\,T_{\mathrm{GT}}/2$.}\label{tab:formats}
\rowcolors{2}{acnmist}{white}
\begin{tabular}{@{}cccccc@{}}
\toprule
\rowcolor{white}Format & Subframes & CP [$\mu$s] & Sequence [$\mu$s] & Guard [$\mu$s] & Range [km] \\
\midrule
0 & 1 & 103.1 & 800  & 96.9  & 14.5 \\
1 & 2 & 684.4 & 800  & 515.6 & 77.3 \\
2 & 2 & 203.1 & 1600 & 196.9 & 29.5 \\
3 & 3 & 684.4 & 1600 & 715.6 & 107.3 \\
\bottomrule
\end{tabular}
\end{table}
Formats with longer CP and guard serve large cells; formats with a repeated sequence collect more
energy for coverage-limited cells. Format~4 (TDD only) uses a short sequence of length 139 in the
special subframe.

\subsection{64 preambles per cell}
Each cell offers 64 preambles. They are generated from one or more ZC roots by cyclic shifts of
$N_{\mathrm{CS}}$ samples: one root yields $\floor{N_{\mathrm{ZC}}/N_{\mathrm{CS}}}$ preambles, and further
roots (consecutive in a broadcast logical order) are used until 64 are obtained. The shift must
exceed the maximum round-trip delay plus delay spread; otherwise a distant UE's preamble would be
detected as a neighbouring shift (\cref{sec:ncs}). The 64 preambles are divided into contention-based
groups (A and B, the latter signalling a larger Msg3) and a set reserved for contention-free use.

\subsection{Contention-based random access (4 steps)}
\begin{enumerate}
  \item \textbf{Msg1 -- preamble.} The UE selects a preamble at random from its group and transmits
        it in the next PRACH occasion, with an open-loop power estimate from the downlink path loss.
  \item \textbf{Msg2 -- random access response (RAR).} Within a response window the eNB sends, on PDSCH
        addressed by an RA-RNTI derived from the PRACH occasion, the detected preamble index (RAPID), a
        timing advance command (11 bits), an uplink grant for Msg3 and a temporary C-RNTI. If no RAR
        arrives, the UE \emph{ramps} its power and retries.
  \item \textbf{Msg3 -- scheduled transmission.} The UE applies the TA and sends, on PUSCH with HARQ, its
        identity and request -- for initial access the \msg{RRC Connection Request} with S-TMSI or a random
        value.
  \item \textbf{Msg4 -- contention resolution.} The eNB echoes the identity it decoded. A UE whose identity
        matches has won; the temporary C-RNTI becomes its C-RNTI. Others start the backoff and retry.
\end{enumerate}
Two UEs that chose the same preamble in the same occasion both receive the same RAR and both send Msg3
on the same resources; typically only one Msg3 is decoded (capture effect), which is why Msg4 must
echo the identity.

\subsection{Contention-free random access}
For handover the target cell reserves a preamble and includes it in the handover command; for
downlink data arrival the eNB signals it on the PDCCH. The UE sends the dedicated preamble, receives
the RAR with TA and grant, and the procedure ends -- Lecture~4 relied on this for the access step of
X2 handover.

\subsection{RACH overload and its control}
When many devices access simultaneously -- after a power outage, at a stadium, or when thousands of
meters report at the same time -- collisions multiply, retries multiply further, and the PRACH collapses.
LTE offers several countermeasures:
\begin{itemize}
  \item \textbf{Access class barring (ACB)}: SIB2 broadcasts a barring factor $a$ and time; before
        accessing, a UE draws a uniform number and proceeds only if it is below $a$, otherwise it waits.
        Extended access barring (EAB) targets delay-tolerant devices.
  \item \textbf{Backoff indicator} in the RAR: losers wait a random time up to the indicated maximum.
  \item \textbf{More PRACH resources}: additional PRACH occasions per frame, at the cost of uplink
        capacity.
  \item \textbf{Separate resources for IoT}: NB-IoT and LTE-M define their own random access channels.
\end{itemize}

\subsection{Random access in 5G NR}
NR keeps the four-step structure but adapts it to beams and to low latency \cite{ts38300}: each
synchronisation signal block (SSB, one per beam) is associated with specific PRACH occasions, so the
preamble also tells the gNB which downlink beam the UE sees best; preambles of length 839 or 139 with
several subcarrier spacings support everything from long-range rural cells to mmWave; and Release~16
adds \textbf{2-step RACH}, in which MsgA combines the preamble and the Msg3 payload and MsgB combines the
RAR and contention resolution, halving the round trips.

% ---------------------------------------------------------------------
\section{Mathematical models}

\subsection{Timing advance}\label{sec:ta}
\paragraph{Problem} How precisely can LTE align a UE, and how far can it be?
\paragraph{Model} The RAR carries $T_A\in\{0,\dots,1282\}$; the UE advances its uplink by
$N_{\mathrm{TA}} = 16\,T_A\,T_\mathrm{s}$ with $T_\mathrm{s} = 1/(15000\times2048)$\,s \cite{ts36213}.
\paragraph{Derivation} One step corresponds to a round-trip time of $16T_\mathrm{s} = \qty{0.521}{\micro\second}$
and therefore to a one-way distance
\begin{equation}\label{eq:tastep}
  \Delta d = \frac{c \cdot 16\,T_\mathrm{s}}{2} = \frac{2.998\times10^8\times0.521\times10^{-6}}{2} \approx \qty{78.1}{\metre},
  \qquad d_{\max} = 1282\,\Delta d \approx \qty{100}{\kilo\metre}.
\end{equation}
\paragraph{Interpretation} The step is a fraction of the normal cyclic prefix (4.7\,$\mu$s), so
residual misalignment is absorbed by the CP. Compare with GSM's 553\,m step (Lecture~1): the finer
step reflects the much shorter LTE symbol and the need to keep all users within the CP. After access,
MAC TA commands (6 bits, relative) keep the alignment as the UE moves.

\subsection{Cyclic shift and cell radius}\label{sec:ncs}
\paragraph{Model} The sequence duration $T_{\mathrm{SEQ}} = 800\,\mu$s spans $N_{\mathrm{ZC}} = 839$ samples,
so one sample lasts $T_{\mathrm{SEQ}}/N_{\mathrm{ZC}} = 0.954\,\mu$s. A UE at distance $r$ shifts its preamble
by the round-trip delay $2r/c$; with multipath delay spread $\tau_{\mathrm{ds}}$ the received energy occupies
$2r/c + \tau_{\mathrm{ds}}$.
\paragraph{Condition} To be unambiguous, this must fit inside one cyclic-shift zone:
\begin{equation}\label{eq:ncs}
  N_{\mathrm{CS}}\,\frac{T_{\mathrm{SEQ}}}{N_{\mathrm{ZC}}} \;\ge\; \frac{2r}{c} + \tau_{\mathrm{ds}}
  \quad\Longrightarrow\quad
  r_{\max} = \frac{c}{2}\left(N_{\mathrm{CS}}\,\frac{T_{\mathrm{SEQ}}}{N_{\mathrm{ZC}}} - \tau_{\mathrm{ds}}\right).
\end{equation}
\paragraph{Example and interpretation} With $\tau_{\mathrm{ds}} = 5.2\,\mu$s: $N_{\mathrm{CS}} = 13$ gives
$r_{\max} \approx \qty{1.1}{\kilo\metre}$ and $\floor{839/13} = 64$ preambles from a \emph{single} root;
$N_{\mathrm{CS}} = 119$ gives $r_{\max}\approx\qty{16.2}{\kilo\metre}$ but only 7 preambles per root, so
$\ceil{64/7} = 10$ roots are needed. Different roots are not orthogonal (cross-correlation
$1/\sqrt{839}\approx0.035$), so large cells have more false alarms and inter-cell preamble interference --
a coverage/robustness trade-off planned per cell. The formula ignores a small detection margin used in
practice; the guard-time limit of \cref{tab:formats} must also hold.

\subsection{Collision probability}\label{sec:collision}
\paragraph{Model} $n$ UEs contend in the same PRACH occasion and choose uniformly among $M$ contention
preambles (e.g.\ $M = 54$ when 10 of the 64 are reserved for CFRA). A UE succeeds if no other UE chose its
preamble (ignoring detection errors and Msg3 capture).
\paragraph{Derivation} The probability that a given UE is alone on its preamble is
\begin{equation}\label{eq:psucc}
  P_\mathrm{s}(n) = \left(1-\frac{1}{M}\right)^{n-1},
\end{equation}
and the expected number of successful UEs in the occasion is
\begin{equation}\label{eq:nsucc}
  S(n) = n\left(1-\frac{1}{M}\right)^{n-1}.
\end{equation}
Treating $n$ as continuous, $\mathrm{d}S/\mathrm{d}n = 0$ gives $n^\ast = -1/\ln(1-1/M)\approx M$, where
$S(n^\ast)\approx M/\e$ -- exactly the slotted-ALOHA efficiency of $1/\e\approx37\%$ per preamble.
\paragraph{Interpretation} For $M = 54$: 10 contenders succeed individually with 84.5\%; 54 contenders give
about 20 successes (37\%); 150 contenders give only 9.3 successes -- \emph{fewer} than 54 did
(\cref{fig:collision}). Beyond the optimum, more load means less throughput, and the failed UEs retry,
pushing the load further: a positive feedback that leads to collapse unless access is throttled (ACB,
backoff). The model ignores power ramping, capture and the multiple PRACH occasions per frame.

\begin{acnfigure}
  \acnfig{lecture-08/collision}
  \caption{Expected successful UEs per PRACH occasion (left axis, \cref{eq:nsucc}) and per-UE success
  probability (right axis, \cref{eq:psucc}) for $M = 54$ contention preambles.}
  \label{fig:collision}
\end{acnfigure}

\subsection{Access delay}
If each attempt succeeds independently with probability $P_\mathrm{s}$, the number of attempts is geometric
with mean $1/P_\mathrm{s}$, and the mean access delay is
\begin{equation}\label{eq:delay}
  \bar D_{\mathrm{RA}} \approx T_{\mathrm{RA}} + \Bigl(\frac{1}{P_\mathrm{s}}-1\Bigr)\bigl(T_{\mathrm{fail}} + \bar T_{\mathrm{BO}}\bigr),
\end{equation}
where $T_{\mathrm{RA}}$ is the duration of a successful 4-step exchange, $T_{\mathrm{fail}}$ the time to
detect a failure (end of the RAR window or of the contention-resolution timer) and $\bar T_{\mathrm{BO}}$ the
mean backoff (\cref{exr:delay}).

% ---------------------------------------------------------------------
\section{Network architecture and diagrams}
Random access is a purely radio-level procedure between UE and eNB (MAC and PHY), but its outcome
feeds the higher layers: Msg3 carries the first RRC message, and Msg4 (in initial access) the
\msg{RRC Connection Setup}, after which the NAS \msg{Service Request} or \msg{Attach Request} reaches the
MME (Lecture~6). \Cref{fig:prach} shows the format-0 preamble, \cref{fig:cbra,fig:cfra,fig:twostep} the
procedures.

\begin{acnfigure}
  \resizebox{\linewidth}{!}{\acnfig{lecture-08/prach-format}}
  \caption{Preamble format 0. The cyclic prefix absorbs the round-trip delay (RTD) and delay spread; the
  guard time keeps a late preamble from overlapping the next subframe.}
  \label{fig:prach}
\end{acnfigure}

% ---------------------------------------------------------------------
\section{Step-by-step procedures}
\begin{procedurebox}{Contention-based random access for initial access}\label{proc:cbra}
\begin{steps}
  \item Read SIB2: PRACH configuration index, root sequence, $N_{\mathrm{CS}}$, preamble groups, power
        parameters, RAR window, contention-resolution timer, ACB parameters.
  \item Pass the access-barring check (if configured), select group A or B, pick a preamble at random and
        transmit it with initial power $P = \min\{P_{\max},\,P_{\mathrm{target}} + \PL\}$.
  \item Monitor the PDCCH for the RA-RNTI during the RAR window. No matching RAR: increase power by the
        ramping step, apply backoff and go back to step~2 (up to the maximum number of attempts).
  \item On RAR: apply TA, store the temporary C-RNTI, send Msg3 on the granted PUSCH, start the
        contention-resolution timer.
  \item On Msg4 with matching identity: success, C-RNTI confirmed, RRC continues. Timer expiry or
        mismatch: backoff and retry.
\end{steps}
\end{procedurebox}

\begin{acnfigure}
  \acnfig{lecture-08/cbra}
  \caption{Contention-based (4-step) random access (\cref{proc:cbra}).}
  \label{fig:cbra}
\end{acnfigure}

\begin{acnfigure}
  \acnfig{lecture-08/cfra}
  \caption{Contention-free random access, e.g.\ at handover.}
  \label{fig:cfra}
\end{acnfigure}

\begin{acnfigure}
  \acnfig{lecture-08/two-step}
  \caption{NR 2-step random access (Release 16).}
  \label{fig:twostep}
\end{acnfigure}

% ---------------------------------------------------------------------
\section{Worked examples}

\begin{examplebox}{Timing advance from a measured delay}\label{ex:ta}
\textbf{Problem.} The eNB measures a preamble round-trip delay of 21.4\,$\mu$s. Which $T_A$ is sent and how
far is the UE?

\textbf{Solution.} $T_A = \mathrm{round}(21.4/0.521) = 41$; $d \approx 41\times78.1 = \qty{3.2}{\kilo\metre}$
(directly: $c\tau/2 = 3.21$\,km).
\end{examplebox}

\begin{examplebox}{Planning $N_{\mathrm{CS}}$}\label{ex:ncs}
\textbf{Problem.} A rural cell must serve UEs up to 12\,km with a delay spread of 5.2\,$\mu$s. Find the
smallest $N_{\mathrm{CS}}$, the preambles per root and the number of roots.

\textbf{Solution.} \Cref{eq:ncs}: $N_{\mathrm{CS}} \ge (2\times12\,000/c + 5.2\times10^{-6})/0.954\times10^{-6} =
(80.1+5.2)/0.954 = 89.4$. The standard set of $N_{\mathrm{CS}}$ values contains 93, which gives
$\floor{839/93} = 9$ preambles per root and $\ceil{64/9} = 8$ roots.
\end{examplebox}

\begin{examplebox}{Event-driven IoT burst}\label{ex:burst}
\textbf{Problem.} After a power cut, 5000 meters in one cell try to access within 1\,s. The cell has one PRACH
occasion per 10\,ms frame (100 per second) and $M = 54$ contention preambles. Estimate the successes if every
meter tries once in a random occasion, and suggest an access-barring factor.

\textbf{Solution.} 50 contenders per occasion on average: $S = 50\,(53/54)^{49} \approx 50\times0.40 = 20.0$
successes per occasion, about 2000 in the first second; the 3000 losers retry and add to the next second's
load. If the 5000 meters instead all fell into 10 occasions (one per 100\,ms), there would be 500 contenders
per occasion and $S = 500\,(53/54)^{499}\approx 0.04$: almost nobody succeeds. A barring factor
$a\approx 54/500\approx0.1$ brings the effective load back to the optimum $n\approx M$, where about $M/\e\approx20$
UEs succeed per occasion.

\textbf{Interpretation.} In an overloaded random access channel, throttling \emph{increases} the throughput --
the central lesson of ALOHA.
\end{examplebox}

% ---------------------------------------------------------------------
\section{Technology comparison}
\begin{table}[htbp]
\centering\small
\caption{Random access across generations.}\label{tab:compare}
\rowcolors{2}{acnmist}{white}
\begin{tabularx}{\linewidth}{@{}l Y Y Y@{}}
\toprule
\rowcolor{white} & GSM & LTE & NR \\
\midrule
Channel           & RACH burst on a timeslot           & PRACH, 6 RBs, ZC preambles        & PRACH tied to SSB beams \\
Contention resolution & SDCCH assignment + identity    & Msg3/Msg4 with identity           & 4-step or 2-step (MsgA/MsgB) \\
TA step / range   & 553\,m / 35\,km                    & 78\,m / 100\,km                   & depends on numerology (finer at higher SCS) \\
Contention-free   & no                                 & yes (handover, DL data)           & yes, plus beam-failure recovery \\
Overload control  & access classes                     & ACB, EAB, backoff                 & unified access control \\
\bottomrule
\end{tabularx}
\end{table}

\paragraph{Trade-offs} Longer preambles and larger cyclic shifts extend range but reduce the number of
orthogonal preambles per root; more PRACH occasions reduce collisions but consume uplink capacity; 2-step
RACH halves latency but needs PUSCH resources for payloads that may collide.

% ---------------------------------------------------------------------
\section{Key takeaways}
\begin{takeawaybox}
\begin{itemize}
  \item Random access provides uplink time alignment and a first identity/grant; it is used for initial
        access, re-establishment, handover and out-of-sync data.
  \item Preambles are ZC sequences ($N_{\mathrm{ZC}} = 839$): constant amplitude, zero cyclic
        autocorrelation; 64 per cell from cyclic shifts of one or more roots.
  \item $r_{\max} = \tfrac{c}{2}(N_{\mathrm{CS}}T_{\mathrm{SEQ}}/N_{\mathrm{ZC}} - \tau_{\mathrm{ds}})$; the guard time also
        bounds the range (format 0: 14.5\,km).
  \item TA step $16T_\mathrm{s} = 0.52\,\mu$s $\Rightarrow$ 78\,m; 11-bit RAR TA $\Rightarrow$ $\approx100$\,km.
  \item CBRA: preamble $\to$ RAR $\to$ Msg3 $\to$ contention resolution; CFRA ends after the RAR.
  \item Successes per occasion $n(1-1/M)^{n-1}$ peak at $\approx M/\e$; overload collapses throughput, so
        throttle (ACB, backoff).
  \item NR adds beam-associated PRACH and 2-step RACH.
\end{itemize}
\end{takeawaybox}

% ---------------------------------------------------------------------
\section{Exercises}

\begin{exercisebox}{1}{Why not just PUSCH?}\label{exr:why}
Explain what would go wrong if an idle UE transmitted its first message directly on PUSCH without
random access. Consider both timing and identity.
\end{exercisebox}

\begin{exercisebox}{1}{CBRA or CFRA?}\label{exr:which}
For each case state whether CBRA or CFRA is used and why: initial attach; handover; RRC re-establishment
after RLF; downlink data arrival for an out-of-sync connected UE.
\end{exercisebox}

\begin{exercisebox}{2}{Format choice}\label{exr:format}
A cell must cover 60\,km over flat terrain. Which LTE preamble formats are possible from the guard-time
point of view, and what is the smallest $N_{\mathrm{CS}}$ (with $\tau_{\mathrm{ds}} = 5.2\,\mu$s)? Is it
possible with an 800\,$\mu$s sequence?
\end{exercisebox}

\begin{exercisebox}{2}{Access delay}\label{exr:delay}
With $T_{\mathrm{RA}} = 20$\,ms, $T_{\mathrm{fail}} = 15$\,ms and a uniform backoff in $[0,20]$\,ms, compute the
mean access delay from \cref{eq:delay} for $n = 10$ and $n = 54$ contenders ($M = 54$).
\end{exercisebox}

\begin{exercisebox}{3}{ZC properties}\label{exr:zc}
Prove that the periodic autocorrelation of \cref{eq:zc} is zero for every non-zero shift $\tau$ when
$N_{\mathrm{ZC}}$ is prime. (Hint: write $x_u(n+\tau)x_u^\ast(n)$ and sum a geometric series in $n$.)
\end{exercisebox}

\begin{exercisebox}{3}{Stability of retries}\label{exr:stab}
New UEs arrive at a mean of $\lambda$ per PRACH occasion; each failed UE retries in the next occasion.
Using \cref{eq:nsucc} with Poisson approximation $S \approx G e^{-G/M}$ for offered load $G$ per occasion,
write the equilibrium condition, find the maximum sustainable $\lambda$ and discuss what happens beyond it.
\end{exercisebox}

\subsection*{Solutions}
\addcontentsline{toc}{subsection}{Solutions}

\begin{solutionbox}{\ref{exr:why}}
Timing: the UE knows only the downlink timing; its uplink would arrive late by the round-trip delay (up to
hundreds of $\mu$s), overlapping other users' SC-FDMA symbols and destroying orthogonality. Identity: the
eNB has no C-RNTI for the UE and cannot address a grant to it; unscheduled PUSCH resources do not exist.
Random access solves both with a robust, delay-tolerant preamble.
\end{solutionbox}

\begin{solutionbox}{\ref{exr:which}}
Initial attach: CBRA (the network does not know the UE). Handover: CFRA (the target reserves a preamble,
saving the contention step and time). RRC re-establishment: CBRA (unplanned). Downlink data for an
out-of-sync UE: CFRA (the eNB knows the UE and signals a dedicated preamble on the PDCCH), falling back to
CBRA if no dedicated preamble is available.
\end{solutionbox}

\begin{solutionbox}{\ref{exr:format}}
The guard time must cover the round trip $2\times60\,000/c = 400\,\mu$s: formats 1 (515.6\,$\mu$s) and 3
(715.6\,$\mu$s) qualify; 0 and 2 do not. \Cref{eq:ncs} requires $N_{\mathrm{CS}} \ge (400+5.2)/0.954 = 425$
samples. The largest non-zero value of the standard unrestricted set is 419 ($r_{\max}\approx59$\,km), just
short; the remaining option is $N_{\mathrm{CS}} = 0$ (no cyclic shift), in which each root provides a single
preamble and 64 roots are used. This works for any delay, but the 64 preambles are no longer orthogonal (only
their low cross-correlation separates them), so detection performance degrades -- the price of very large
cells. The 800-$\mu$s sequence length is not the issue; repetition (format 3) helps the link budget only.
\end{solutionbox}

\begin{solutionbox}{\ref{exr:delay}}
$\bar T_{\mathrm{BO}} = 10$\,ms. $n=10$: $P_\mathrm{s} = 0.845$, $1/P_\mathrm{s}-1 = 0.183$, $\bar D = 20 + 0.183\times25 =
24.6$\,ms. $n=54$: $P_\mathrm{s} = 0.371$, $1/P_\mathrm{s}-1 = 1.70$, $\bar D = 20 + 1.70\times25 = 62.4$\,ms
(optimistic, because retries also add load).
\end{solutionbox}

\begin{solutionbox}{\ref{exr:zc}}
$x_u(n+\tau)x_u^\ast(n) = \exp(-\jj\pi u[(n+\tau)(n+\tau+1) - n(n+1)]/N) =
\exp(-\jj\pi u[\tau^2+\tau]/N)\,\exp(-\jj 2\pi u\tau n/N)$. The first factor is constant in $n$; summing the
second over $n = 0,\dots,N-1$ gives $\sum_n \e^{-\jj2\pi u\tau n/N} = 0$ unless $u\tau \equiv 0 \pmod N$. With
$N$ prime and $1\le u,\tau<N$, $u\tau$ is never a multiple of $N$, so the autocorrelation is zero for all
$\tau\neq0$. (Periodicity of the quadratic phase modulo $N$ holds for odd $N$.)
\end{solutionbox}

\begin{solutionbox}{\ref{exr:stab}}
In equilibrium the success rate equals the arrival rate: $\lambda = G e^{-G/M}$. The right-hand side is
maximal at $G = M$ with value $M/\e$, so $\lambda_{\max} = M/\e \approx 19.9$ UEs per occasion for $M=54$. For
$\lambda < M/\e$ there are two equilibria: the low-load one ($G<M$) is stable, the high-load one is not. For
$\lambda > M/\e$ no equilibrium exists: the backlog grows without bound and successes fall towards zero --
the instability that ACB and randomised backoff are designed to prevent.
\end{solutionbox}

% ---------------------------------------------------------------------
\section{Further reading}
The MAC random access procedure is in \cite{ts36321}, preamble generation and formats in \cite{ts36211},
timing advance in \cite{ts36213}, and ACB in the RRC specification \cite{ts36331}. NR random access,
including 2-step RACH, is summarised in \cite{ts38300}. Sesia et al.\ \cite{sesia2011lte} contains a detailed
chapter on PRACH design and ZC sequences; Dahlman et al.\ \cite{dahlman2020nr} covers the NR procedures.

\printbibliography[heading=subbibliography,title={\ACNstrRefs}]
\end{document}
